% Part: computability
% Chapter: computability-theory
% Section: russells-paradox

\documentclass[../../../include/open-logic-section]{subfiles}

\begin{document}

\olfileid{cmp}{thy}{rus}
\olsection{Bandhingan karo Paradoks Russell}

Migunani yen argumen-argumen ing bagean iki
dibandhingake karo paradoks Russell, kalebu
padha lan bedane:
\begin{enumerate}
\item Paradoks Russell: tetepna $S = \Setabs{x}{x \notin x}$.
  % OLPL-136: Sumber beku nulis X ora kalebu S; substitusi S ing pambentukan S mbutuhake S ora kalebu S.
  Mula $S \in S$ yen lan mung yen $S \notin S$,
  sawijining kontradiksi. Cathetan koreksi sumber OLPL-136: X kang ora ditegesi ing rumus Inggris dibenerake dadi S, manut substitusi S ing pambentukan himpunan kasebut.

  \emph{Dudutan:} Himpunan~$S$ kaya mangkono ora ana.
  Panganggep yen ana ``himpunan kabeh himpunan''
  nalisir aksioma liyane ing teori himpunan.

\item Owah-owahan paradoks Russell: tetepna $F$ minangka
  ``fungsi'' saka himpunan kabeh fungsi menyang
  $\{ 0, 1 \}$, kang ditegesi minangka
  \[
  F(f) =
  \begin{cases}
    1 & \text{yen $f$ kalebu domain $f$, lan $f(f) = 0$} \\
    0 & \text{yen ora}
  \end{cases}
  \]
  Argumen kang padha nuduhake yen $F(F) = 0$
  yen lan mung yen $F(F) = 1$, sawijining kontradiksi.

  \emph{Dudutan:} $F$ dudu fungsi. ``Himpunan kabeh
  fungsi'' kegedhen kanggo dadi domain sawijining fungsi.

\item Argumen diagonalisasi: tetepna $f_0$, $f_1$, \dots
  minangka enumerasi fungsi komputabel parsial,
  lan tetepna $G \colon \Nat \to
  \{ 0, 1 \}$ ditegesi minangka
  \[
  G(x) =
  \begin{cases}
    1 & \text{yen $f_x(x)\downarrow = 0$} \\
    0 & \text{yen ora}
  \end{cases}
  \]
  Yen $G$ komputabel, mula iku fungsi $f_k$
  kanggo sawenehing $k$. Nanging banjur $G(k) = 1$
  yen lan mung yen $G(k) = 0$, sawijining kontradiksi.

  \emph{Dudutan:} $G$ ora komputabel. Gatekna yen miturut
  aksioma teori himpunan, $G$ tetep sawijining fungsi;
  ing kene ora ana paradoks, mung ana panjlentreh.
\end{enumerate}

Rembugan bab fungsi parsial, fungsi komputabel,
fungsi komputabel parsial, lan sapanunggalane
bisa mbingungake. Himpunan kabeh fungsi parsial
saka $\Nat$ menyang $\Nat$ minangka kumpulan
obyek kang gedhe. Sawetara fungsi kasebut total,
sawetara komputabel, sawetara total lan komputabel,
lan sawetara dudu loro-lorone. Elinga yen nalika
kita nyebut ``fungsi'' tanpa katrangan liyane,
kang dimaksud lumrahe fungsi total. Mula ana:
\begin{enumerate}
\item fungsi komputabel
\item fungsi komputabel parsial kang ora total
\item fungsi kang ora komputabel
\item fungsi parsial kang ora total lan ora komputabel
\end{enumerate}
Kanggo mbedakake iki kabeh, bisa migunani
yen kokgambar kothak gedhe kang nggambarake
kabeh fungsi parsial saka $\Nat$ menyang $\Nat$,
banjur tandhana rong wilayah kang tumpang tindih,
yaiku wilayah fungsi total lan wilayah fungsi
komputabel parsial. Coba terangna conto sawijining
obyek ing saben wilayah kang kawangun ing gambar iku.

\end{document}
